\section{DGLAP 演化结构}

\subsection{纯胶动力学中 DGLAP 何时为对角}

由式 (\ref{M03}) 定义的组合 $M_{03} ^+\equiv
{ M}_{0 i; i 3} + { M}_{3 i;i 0}$ 含有扭度 2 振幅 $\mathcal{M}_{pp}$：
\begin{align}
M_{03} ^+
= 4 p_0 p_3 \mathcal{M}_{pp} +
2 p_0 z_3 \left ( \mathcal{M}_{pz} +\mathcal{M}_{zp} \right) \ ,
\end{align}
但混有更高扭度的成分 \mbox{$\mathcal{M}_{zp}+ \mathcal{M}_{pz}$}。
在定域极限下，相关算符正比于坡印廷矢量的第三分量
\begin{align*}
{\bf S_3}=
({\bf E \times B})_3 =E_1 B_2- B_1 E_2 =-( G_{01} G_{13} +G_{32} G_{20}) \ .
\end{align*}

如前所述，纯胶动力学中矩阵元
$M_{03} ^+(z_3,p)$ 单圈修正的盒图部分\footnote{这种简单性在高阶可能被破坏。}具有简单的 DGLAP 结构 (\ref{box+})。
把所有单圈胶子修正合并起来，
在 $\overline{\rm MS}$ 方案中我们得到
\begin{align}
& \frac{g^2N_c}{8\pi^2}\int_0^1 \dd u
\left\{ \left[ \left (\frac{3}{2} -\frac16 \right ) \ln(z_3^2 \mu_{\rm UV}^2 e^{2\gamma_E}/4) +2 \right] \delta(\bar u) \right. \nonumber \\
&\quad \left. - 2 \log(z_3^2 \mu_{\rm IR}^2 e^{2\gamma_E}/4) \left [ \frac{(1- u \bar u)^2}{\bar u} \right ]_+ \right.
\label{M031L} \\
&\quad \left. + \left [ u -3 \frac{u}{\bar u} -4 \frac{\log (\bar u)}{\bar u} \right ]_+
+ 2 \left (\bar uu + \frac{2}{3}\bar u^3 \right ) \right\} \ M_{03} ^+(uz,p) \ .\nonumber
\end{align}

第一行来自 UV 奇异贡献。它含 $\delta(\bar u)$ 因子，
反映 UV 发散的定域性质，并把
$M_{03} ^+(uz,p)$ 变为 $M_{03} ^+(z,p)$。
第二行含 Altarelli-Parisi（AP）核
\begin{align}
B_{gg} (u) =&
2 \left [\frac{(1-u\bar u)^2} {1-u} \right ]_{+} \ .
\label{V1}
\end{align}
其加 prescription 结构反映了如下事实：
在定域极限下 ${\mathcal{M}}_{pp}(z,p)$ 正比于
胶子能量动量张量的矩阵元。从现在起，``+''表示在 1 处的加 prescription。

第三行
含来自顶点图（它们具有加 prescription 形式）与盒图的与 $z_3$ 无关的项。
后者可写成具有加 prescription 形式的项 $2 \left (\bar uu +  \frac{2}{3}\bar u^3  \right )_{+}$
与可并入 UV 项的项 $\frac23 \delta (\bar u)$ 之和。

\subsection{约化 Ioffe 时间分布}

组合 $\mathcal{M}_{pz} +\mathcal{M}_{zp}$ 是 $\nu =z_3 p_3$ 的奇函数。
把它写作 $2 z_3 p_3 m^+_{zp} (\nu, z_3^2)$，便有
\begin{align}
M_{03} ^+(z_3, p) = 4 p_3 p_0 [ \mathcal{M}_{pp} (\nu, z_3^2 ) +
z_3^2 \ m^+_{zp} (\nu, z_3^2 )] \ .
\label{M03p}
\end{align}
除掉运动学因子 $4 p_3 p_0$ 后，我们处理的是
\begin{align}
\widetilde{\mathcal{M}}_{pp}(\nu, z_3^2 ) \equiv
\mathcal{M}_{pp}(\nu, z_3^2 ) +
z_3^2 \ m^+_{zp} (\nu, z_3^2 ) \ ,
\end{align}
它是 $\nu$ 与 $z_3^2$ 的函数。
现在，正如文献 \cite{Radyushkin:2017cyf,Orginos:2017kos} 中考虑的夸克情形那样，我们可以引入
约化 Ioffe 时间分布
\begin{align}
\widetilde{\mathfrak M} (\nu, z_3^2) \equiv \frac{\widetilde{\mathcal{M}}_{pp}(\nu, z_3^2)}{\widetilde{\mathcal{M}}_{pp} (0, z_3^2)} \ .
\label{redm0}
\end{align}
由于 $\widetilde{\mathcal{M}}_{pp}(\nu, z_3^2)$ 来自乘法可重正的组合 $M_{03} ^+$，
链接相关图与胶子自能图产生的
UV 发散因子 $Z(z_3^2 \mu^2_{UV})$ 在比值 (\ref{redm0}) 中相消。于是，约化
伪 ITD $\widetilde{\mathfrak M} (\nu, z_3^2)$ 的小 $z_3^2$ 依赖仅来自对数型 DGLAP 演化
效应。
此外，由于 AP 核 $B_{gg} (u)$ 的加 prescription 性质，
$\widetilde{\mathcal{M}}_{pp}(0, z_3^2)$
没有 DGLAP 型的对数 $z_3^2$ 依赖。

于是，略去 ${\cal O} (z_3^2)$ 项，我们得出结论：{\it 在纯胶动力学中，}
$\widetilde{\mathfrak M} (\nu, z_3^2)$ 满足关于
$z_3^2$ 的演化方程
\begin{align}
\frac{d}{d \ln z_3^2} \,
\widetilde{\mathfrak M} (\nu, z_3^2) &= - \frac{\alpha_s}{2\pi} \, N_c
\int_0^1 du \, B_{gg} (u) \widetilde{\mathfrak M} (u \nu, z_3^2)
\label{EE}
\end{align}
当存在胶子--夸克跃迁时该关系将被修正。

\subsection{胶子--夸克混合}

在 $\overline{\rm MS}$ 方案中，图 \ref{gluqum} 所示胶子--夸克图对 $M_{03} ^+$ 的贡献为
\begin{align}
& \frac{g^2 C_F}{4\pi^2 z_3} \int_0^1 \dd u
\Biggl [ -2u - \ln(z_3^2 \mu_{\rm IR}^2 e^{2\gamma_E}/4) \left[2\bar u+\delta(\bar u) \right] \Biggr ] {\cal O}_q (uz_3) \ ,
\label{gqmix}
\end{align}
其中 ${\cal O}_q (z_3)$ 是夸克场的单态组合：
\begin{align}
& {\cal O}_q (z_3)= \frac{i}{2} \sum_f \left(\bar\psi_f( 0) \gamma_0 \psi_f (z_3) - \bar\psi_f( z_3) \gamma_0 \psi_f(0) \right) \ ,
\label{gsing}
\end{align}
$f$ 标记夸克味。
注意 ${\cal O}_q (z_3)$ 在 $z_3=0$ 时为零。
把 ${\cal O}_q (z_3)$ 按 $z_3$
展开：
\begin{align}
& {\cal O}_q ( z_3)= z_3\, \frac{i}{2} \sum_f \bar\psi_f( 0)\gamma_0 \!\stackrel{\leftrightarrow}{\partial}_3\! \psi_f(0) + {\cal O} (z_3^3) \ ,
\label{gsingb}
\end{align}
可见最低阶项正比于能量动量张量的夸克部分。
该项带有 $z_3$ 因子，恰好抵消式 (\ref{gqmix}) 中整体的 $1/z_3$ 因子。

${\cal O}_q ( z_3)$ 的矩阵元可参数化为
\begin{align}
& \langle p| {\cal O}_q ( z_3) |p\rangle = 2 p^0
\int_0^1 \dd x \, \sin (x p_3 z_3) \, q_S (x)
\label{qpar}
\end{align}
其中 $f_S (x) = \sum_f [q_f (x) +\bar q_f (x)]$ 是单态夸克分布。
为提出整体的 $z_3$ 因子，我们改写为
\begin{align}
& \langle p| {\cal O}_q ( z_3) |p\rangle
= 2 p^0 p_3 z_3
\int_0^1 \dd y \, \int_0^1 \dd \alpha\, \cos (\alpha y \nu) \, y \, f_S (y )
\label{qpar2}
\end{align}
其中照例 $\nu = p_3 z_3$。这给出
\begin{align}
&\frac1{z_3} \int_0^1 \dd u \, A(u) \, \langle p| {\cal O}_q (u z_3) |p\rangle
\nn &
= 2 p^0 p_3 \, \int_0^1 \dd w {\cal I}_S (w \nu) {\cal A} (w)
\label{qpar3}
\end{align}
适用于式 (\ref{gqmix}) 类型的 $u$ 积分。此处
\begin{align}
& {\cal I}_S (\nu) =
\int_0^1 \dd y \ \cos (y \nu) \, y\, f_S (y)
\label{MS}
\end{align}
是单态夸克 Ioffe 时间分布，而
\begin{align}
{\cal A} (w)=
\int_w^1 \dd u \, A(u) \ .
\label{Aw}
\end{align}
对演化核 $B_{gq} (u) \equiv 2\bar u+\delta(\bar u)$，我们得到
\begin{align}
{\cal B}_{gq} (w)=
\int_w^1 \dd u \, B_{gq} (u) = 1+ (1- w)^2 \ .
\label{Bw}
\end{align}

\subsection{匹配关系}

$M_{03} ^+(z_3, p)$ 组合的一个缺点是
它在 $p_3=0$ 时为零（见 \mbox{式 (\ref{M03p})）}。
因此要提取 $\nu=0$ 的
$\widetilde{\mathcal{M}}_{pp}(\nu, z_3^2 )$，
需要在若干较小的 $p_3$ 值处测量 $M_{03} ^+(z_3, p)$，
除去 $p_3$，
并把结果外推到 \mbox{$p_3=0$}。这一手续会带来额外的系统不确定度。

\begin{figure}[t]
\centerline{\includegraphics[width=1.7in]{images/gluqum}}
\vspace{-3mm}
\caption{胶子--夸克混合图。\label{gluqum}}
\end{figure}

幸运的是，
\mbox{式 (\ref{00m}),} 的组合
${ M}_{0 i; i 0} - { M}_{i j; i j} = 2 p_0^2 \mathcal{M}_{pp}$
正比于 $p_0^2$，不存在这个问题。
而且它给出无污染的扭度 2 振幅 $\mathcal{M}_{pp}$。
如此得到的振幅 $\mathcal{M}_{pp} (\nu, z_3^2)$
可用来按式 (\ref{redm0}) 构造约化伪 ITD
${\mathfrak M} (\nu, z_3^2)$。

利用我们对
${ M}_{0 i; i 0}$ 与 ${ M}_{i j; i j}$ 单圈修正的计算结果，并在修正中只保留
${\cal M}_{pp}$ 项（略去``高扭度''项
${\cal M}_{zz}, {\cal M}_{zp}, {\cal M}_{pz}, {\cal M}_{ppzz}$），
我们得到匹配关系
\begin{align}
{\mathfrak M}& (\nu, z_3^2)\, {{\cal I}_g (0, \mu^2) } = { {\cal I}_g (\nu, \mu^2) }
- \frac{\alpha_s N_c}{2\pi}\int_0^1 \dd u \, {{\cal I}_g (u \nu, \mu^2)}
\nonumber \\
& \times
\Biggl \{ \ln (z_3^2 \mu^2 e^{2\gamma_E}/4) \, B_{gg} (u) +4 \left [ \frac{u+ \log (\bar u)}{\bar u}\right ]_{+} \nonumber \\
& + \frac23 \left [ 1 -u^3
\right ]_+ \Biggr \}
- \frac{\alpha_s C_F}{2\pi} \ln (z_3^2 \mu^2 e^{2\gamma_E}/4)
%
\nonumber \\
&
\times
\int_0^1 \dd w \,
\left [{\cal I}_S (w \nu, \mu^2) - {\cal I}_S (0, \mu^2) \right ]
\, {\cal B}_{gq} (w) \
\label{matching}
\end{align}
它联系``格点函数''
${\mathfrak M}( \nu,z_3^2)$ 与光锥 ITD ${\cal I}_g (\nu,\mu^2)$ 及
${\cal I}_S(\nu,\mu^2)$。前者通过
\begin{align}
{\cal I}_g (\nu,\mu^2) =\frac12
\int_{-1}^1 dx \, \,
e^{ix\nu} \,x {f}_g(x,\mu^2) \ .
\label{If}
\end{align}
与胶子 PDF ${f}_g(x,\mu^2)$ 相联系。
由于 $x {f}_g(x,\mu^2)$ 是 $x$ 的偶函数，
${\cal I}_g (\nu,\mu^2)$ 的实部由 $x {f}_g(x,\mu^2)$ 的余弦变换给出，
虚部为零。
因子 ${\cal I}_g (0, \mu^2)$ 的含义是
胶子携带的强子动量份额，${\cal I}_g (0, \mu^2) = \langle x \rangle_{\mu^2}$。

于是，式 (\ref{matching}) 可用来提取胶子分布的形状。
其归一化，即 $\langle x \rangle_{\mu^2}$ 的值，应当由独立的格点计算确定，
类似 Ref.~\cite{Yang:2018bft} 中所做的那样。
出现在 ${\cal O} (\alpha_s)$ 修正中的单态夸克函数 ${\cal I}_S (w \nu, \mu^2)$
也应当独立地计算（或估计）。

把式 (\ref{If}) 代入匹配条件 (\ref{matching})，我们可将后者改写为核形式 \cite{Radyushkin:2019owq}
\begin{align}
{\mathfrak M}(\nu,z_3^2) & =
\int_{0}^1 \dd x \,
\frac{x {f}_g (x,\mu^2) }{\langle x \rangle_{\mu^2}} \, R_{gg} (x \nu, z_3^2 \mu^2 ) \nn &
+
\int_{0}^1 \dd x \,
\frac{x {f}_S (x,\mu^2) }{\langle x \rangle_{\mu^2}} \, R_{gq} (x \nu, z_3^2 \mu^2 ) \
,
\label{MtchI}
\end{align}
其中核 $R_{gg} (x \nu, z_3^2 \mu^2 )$
由下式给出：
\begin{align}
R_{gg}(y, z_3^2 \mu^2) =& \cos y
- \frac{\alpha_s}{2\pi} \, N_c
\Biggl \{ \ln \left ( z_3^2 \mu^2\frac{e^{2\gamma_E+2}}{4} \right ) \,
R_B(y)\nn & +R_U (y)
+R_L(y) +R_C(y) \Biggr \} \, ,
\label{MtchR}
\end{align}
其中 $R_B(y)$ 是
$B_{gg}$ 核的余弦傅里叶变换：
\begin{align}
R_B(y) =& \int_0^1 \dd u \, B_{gg} (u) \, \cos (u y) \ .
\label{RB}
\end{align}
其计算直截了当，结果可用 $\cos y$、$\sin y$
以及积分余弦 $\text{Ci}(y)$ 和积分正弦 $\text{Si}(y)$ 函数表示。
后者来自 $B(u)$ 的 $1/(1-u)$ 部分，该部分给出
\begin{align}
\int_0^1 \dd u \, & \left[ \frac1{1-u} \right ]_+ \, \cos (u y) = \sin (y) \, \text{Si}(y) \nn & +\cos (y)\,
[\text{Ci}(y)-\log (y)-\gamma_E ]
\ .
\label{RB0}
\end{align}
此组合也出现在 $R_U(y)$ 中，后者是 $4[u/\bar u]_+$ 的余弦变换。
类似地，$R_L(y)$ 是 \mbox{$4[(\ln \bar u)/\bar u]_+$} 项的余弦变换，
它涉及一个超几何函数
\begin{align}
R_L(y) =
4\, {\rm Re} \left [ i y e^{i y} \,
_3F_3(1,1,1;2,2,2;-i y) \right ] \, .
\label{MtchRL}
\end{align}

核 $R_C(y)$ 与 $R_{gq} (y)$ 分别由 $\frac23 \left [ 1-u^3
\right ]_+$ 与 $\left [ 1+\bar u^2\right]_+$ 的余弦变换给出。
它们的表达式
涉及 $\cos y$、$\sin y$ 及 $y$ 的负幂。

关键的一点是：$R(y, z_3^2 \mu^2)$ 核由{\it 显式的}可微扰计算的表达式给出。利用它们以及式 (\ref{MtchI})，可以直接联系
${\mathfrak M}(\nu,z_3^2)$ 与光锥
PDF。然后，
对 $f_g(x,\mu^2)$ 与 $f_S(x, \mu^2)$ 分布假设某些参数化形式，
便可利用式 (\ref{MtchI})、(\ref{MtchR}) 从 ${\mathfrak M}(\nu,z_3^2)$ 的格点数据
拟合其参数与 $\alpha_s$。这一步骤本质上与
``好格点截面''方法 \cite{Ma:2014jla,Ma:2017pxb} 所用的相同。

\subsection{准分布的匹配关系}

核关系 (\ref{MtchI}) 直接把 ${\mathfrak M}(\nu,z_3^2)$
与 PDF 联系起来。
因此无需引入
准分布之类的中间函数。
尽管如此，我们对特定矩阵元的结果，
例如 $M_{03} ^+(z_3,p)$ 的式 (\ref{M031L})，
仍可用于求准分布的匹配条件。后者一般定义为 \cite{Ji:2013dva}
\begin{align}
Q(y, p_3) = \frac{p_3}{2 \pi}
\int_{-\infty}^\infty \dd z_3
{\cal M} (z_3, p) \, e^{-i y p_3 z_3} \
.
\label{qPDF}
\end{align}

为此，应把振幅 ${\cal M} (z_3, p)$ 通过核
关系 (\ref{MtchI}) 写出，并把 $R(x \nu, z_3^2 \mu^2 )$ 用 $p_3$ 与 $z_3$ 表达，
记作 $J(x, p_3, z_3)$。
其在单圈对 $z_3$ 依赖的结构可从式 (\ref{M031L}) 读出。
对 $gg$ 部分，
\begin{align}
& J_1^{gg} (x, p_3, z_3)= \left (\gamma_U \ln z_3^2 +C_U \right) e^{ix p_3 z_3}
\nn & - \int_0^1 \dd u \,\left [ \ln z_3^2 \, B_{gg} (u)
+C (u) \right ]e^{iu x p_3 z_3} \ .
\label{Jk}
\end{align}
单圈准 PDF 匹配核随即由下式给出：
\begin{align}
Z_1^{gg} (y, x, p_3) = \frac{p_3}{2 \pi}
\int_{-\infty}^\infty \dd z_3 \, J_1^{gg} (x, p_3, z_3)
\, e^{-i y p_3 z_3} \
.
\label{qZ}
\end{align}
$J_1^{gg} (x, p_3, z_3)$ 中 $C_U$ 与 $C(u)$ 的贡献
产生 \mbox{$C_U \delta (y- x)$} 与 $C (u)\delta (y - u x )$ 项。
因此，
$Q(y,p_3)$ 的这些部分只在``正则''区域 $0\leq |y| \leq 1$ 内可见。然而，
含 $\ln z_3^2$ 的项在 $y/x>1$
与 $y/x<0$ 区域也给出非零贡献，即
\begin{align}
& \, Z_1^{gg} (y, x, p_3) |_{y/x>1} = \frac1{|x|} \left [- \frac{\gamma_U}{\eta -1}
+ \int_0^1 \dd u \, \frac{B_{gg} (u)}{\eta-u} \right ]
\ ,
\label{Zy}
\end{align}
其中 $\eta= y/x$。
$Z_1^{gg} (y, x, p_3) |_{y/x<0}$ 项由同一表达式冠以相反符号给出。
注意这些贡献完全
由 AP 核 $B_{gg} (u)$ 与 UV 常数
$\gamma_U$ 决定。
知道它们之后，可由式 (\ref{Zy}) 导出对任何费曼图计算所得的
$Z_1^{gg} (\eta, 1, p_3) |_{\eta>1}$
与 $Z_1^{gg} (\eta, 1, p_3) |_{\eta<0}$
结果的一般约束。
利用 $B_{gg} (u)$ 的显式形式，我们求得
\begin{align}
\int_0^1 \dd u \, \frac{B_{gg} (u) }{\eta-u} & = 2 \frac{(1-\eta \bar \eta)^2}{\eta -1} \, \ln \frac{\eta-1}{\eta}
\nn & + \frac{11}{6} \frac{1}{\eta -1} + \eta (2 \eta -1) +\frac{11}{3}
\ .
\label{ZB}
\end{align}

对大的 $\eta$，此表达式以 ${\cal O}(1/\eta^2)$ 趋于零。
应当强调，任何具有加 prescription 形式的核 $B(u)$ 都会导致这种行为。

这一观察连同 \mbox{式 (\ref{ZB})}
给出的显式表达式可用于检验
\mbox{文献 \cite{Wang:2017qyg,Wang:2019tgg}} 中的胶子--胶子匹配核。
我们的检验表明，对应于 Ref.~\cite{Wang:2019tgg}
式 (64) 的 $Z_1^{gg} (y, 1, p_3) |_{y>1}$ 不满足约束 (\ref{Zy})。
差别是一个常数项 (-2/3)，它导致
$Z_1^{gg} (y, 1, p_3) |_{y>1}$ 对 \mbox{$y$} 的积分线性发散。
同样的差别（符号相反）出现在 Ref.~\cite{Wang:2019tgg} 式 (64) 的
$Z_1^{gg} (y, 1, p_3) |_{y<0}$ 项中。
显然，这些差别源于 Ref.~\cite{Wang:2019tgg} 计算中使用了离壳外胶子场，
但对这一话题的讨论超出了本文的范围。
